\documentclass[a4paper,12pt]{article}

% Pacotes essenciais
\usepackage[utf8]{inputenc}
\usepackage[T1]{fontenc}
\usepackage[brazil]{babel}
\usepackage{amsmath, amssymb}
\usepackage{graphicx}
\usepackage{geometry}
\usepackage{hyperref}
\usepackage{float}
\usepackage{indentfirst}
\usepackage{setspace}

\onehalfspacing

\geometry{a4paper, margin=2.5cm}

\begin{document}

\begin{titlepage}
    \centering

    UNIVERSIDADE FEDERAL DE SERGIPE\\
    CIDADE UNIVERSITÁRIA PROFESSOR JOSÉ ALOÍSIO DE CAMPOS

    \vspace{4cm}

    FRANCISCO PASSOS DOS SANTOS ALVES

    \vspace{4cm}

    NOTAS DE ESTUDO:\\
    FUNDAMENTOS TEÓRICOS DA CIÊNCIA DA COMPUTAÇÃO
    \vfill

    SÃO CRISTÓVÃO\\
    2026

\end{titlepage}

\pagenumbering{arabic}

\tableofcontents

\newpage
\section{Introdução}

Desde o átimo em que pisei na universidade, não me restavam dúvidas sobre meu desejo de seguir  a carreira acadêmica, exercendo o ofício não apenas de 
pesquisador, mas também de professor e divulgador científico. Estas notas sobre os fundamentos teóricos da ciência da computação reúnem tudo aquilo que 
mais respeito dentro da computação.

O principal objetivo deste documento é reunir minhas anotações, produzir novas e deixá-las disponíveis publicamente para graduandos que apreciam matemática 
e computação teórica tanto quanto eu, ou para aqueles que desejam explorar tais assuntos com um pouco mais de profundidade, sem a pressa \textit{necessariamente} 
imediatista e ridícula imposta sobre os alunos, e sem a burocracia à qual estão submetidos.

Todo o conhecimento aqui presente será referenciado ao seu respectivo \textit{habitat} de origem, seja em livros, artigos, documentações, notas de algum  professor 
obtidas com algum aluno aleatório da Universidade Federal de Sergipe (UFS), ou mesmo materiais encontrado na internet. Serão todos devidamente referenciados.

Este documento é disponibilizado e atualizado oficialmente no repositório \url{https://github.com/FrankSteps/formal-languages-notes}. Sugestões de conteúdo ou 
correções são bem-vindas: sinta-se livre para abrir uma \textit{issue} ou enviar um email para \href{mailto:franciscopassos.contato@gmail.com}{franciscopassos.contato@gmail.com}. 
Reze para que eu não tenha mudado o \textit{username} ou o nome do repositório.

Aproveite o conteúdo.

\newpage

\section{Fundamentos Matemáticos da Computação}

\subsection{Noções e terminologias matemáticas}

\subsubsection{Conjuntos} % sipser + fundamentos elementares da matemática: conjuntos
\subsubsection{Sequências e uplas} % guidorizzi
\subsubsection{Funções e relações} % widman
\subsubsection{Grafos} % sipser + curso da univesp
\subsubsection{Cadeias e linguagens} % sipser


\subsection{Métodos de Provas} % curso da univesp

\subsubsection{Prova direta}
\subsubsection{Prova por contraposição}
\subsubsection{Prova por contradição}
\subsubsection{Prova por indução}
\subsubsection{Prova por divisão de casos}

\subsection{Fundamentos da teoria dos grafos}


\section{Linguagens Formais e Autômatos}

\subsection{Autômato finito determinístico}
\subsection{Autômato finito não determinístico}

\section{Máquinas de Estado: Moore e Mealy}

\section{Computabilidade e Máquinas de Turing}

\section{Cálculo Lambda}

\section{Teoria dos Tipos}

\section{Estruturas Algébricas}

\section{Teoria das Categorias}

\section{Considerações Finais}

\newpage
\section{Referências}


\begin{thebibliography}{99}

\bibitem{widman}
WIDMAN, Jack. \textit{Aprenda Programação Funcional}. São Paulo: Alta Books, 2025.

% \bibitem{lipovaca}
% LIPOVAČA, Miran. \textit{Learn You a Haskell for Great Good!}. San Francisco: No Starch Press, 2011.

\bibitem{rojas}
ROJAS, Raul. \textit{A tutorial introduction to the lambda calculus}. Versão 2.0. Berlin: Freie Universität Berlin, 2015.

\bibitem{sipser}
SIPSER, Michael. \textit{Introdução à Teoria da Computação}. Tradução da 2ª edição norte-americana. São Paulo: Cengage Learning, 2007.

\bibitem{pinter}
PINTER, Charles C. \textit{A Book of Abstract Algebra}. 2nd ed. Mineola, N.Y.: Dover Publications, Inc., 2010.

\bibitem{univesp}
POSSANI, Claudio. \textit{Fundamentos Matemáticos da Computação}. Univesp, YouTube, 2017. Disponível em: \url{https://www.youtube.com/playlist?list=PLxI8Can9yAHcXBgFryV0AV7LYdLR1skuF}. Acesso em: 8 de agosto. 2026.

\end{thebibliography}
\end{document}